% !TEX program = lualatex
% Econometrica manuscript on the Econometric Society class (econsocart.cls,
% econsocart.cfg, ecta-fullname.bst). Option [draft] typesets the submission
% version; [final] the prepublication version. Grey angle-bracket slots mark
% the content each sentence carries.
\documentclass[ecta,nameyear,draft]{econsocart}

\RequirePackage[colorlinks,citecolor=blue,linkcolor=blue,urlcolor=blue,pagebackref]{hyperref}
\usepackage{luatexja-fontspec}
\usepackage{xcolor}
\usepackage{amssymb}
\usepackage{bm}
\usepackage{float}
\usepackage{booktabs}
\usepackage{longtable}
\usepackage{threeparttable}
\ifdefined\abovecaptionskip\else\newlength{\abovecaptionskip}\fi
\usepackage{graphicx}
\newcommand{\hiddenlabel}[1]{\refstepcounter{equation}\label{#1}}
\newcommand{\hiddenfiglabel}[1]{\refstepcounter{figure}\label{#1}}
\newcommand{\hiddentablabel}[1]{\refstepcounter{table}\label{#1}}

\startlocaldefs
\newcommand{\econtheorem}[3]{\ifcsname #1\endcsname\else\theoremstyle{#2}\newtheorem{#1}{#3}\fi}
\econtheorem{assumption}{plain}{Assumption}
\econtheorem{proposition}{plain}{Proposition}
\econtheorem{lemma}{plain}{Lemma}
\econtheorem{corollary}{plain}{Corollary}
\econtheorem{definition}{remark}{Definition}
\econtheorem{example}{remark}{Example}
\ifdefined\slot\else\DeclareRobustCommand{\slot}[1]{{\color{gray}\itshape\ensuremath{\langle}#1\ensuremath{\rangle}}}\fi
\newcommand{\E}{\mathbb{E}}
\newcommand{\Var}{\operatorname{Var}}
\endlocaldefs

\begin{document}

\begin{frontmatter}

%% BEGIN TITLE
\title{\slot{Title}}
\runtitle{\slot{Running title}}
%% END TITLE

\begin{aug}
%% BEGIN AUTHORS
\author[id=au1,addressref={add1}]{\fnms{First}~\snm{Author}\ead[label=e1]{author@university.edu}}
\address[id=add1]{\orgdiv{Department of Economics}, \orgname{University}}
%% END AUTHORS

%% BEGIN SUPPORT
\support{We thank \slot{research assistants} for research support and \slot{discussants}, the editor, three anonymous referees and seminar participants for comments. \slot{Author} received financial support from \slot{funder}.}
%% END SUPPORT

%% BEGIN ABSTRACT
\begin{abstract}
This paper studies the effect of \slot{treatment} on \slot{outcome}. We focus on \slot{setting}, which \slot{scale and reach}. Using \slot{data, population and period}, we estimate \slot{estimand}. \slot{Source of variation} assigns \slot{treatment} independently of \slot{potential outcomes}. We find that \slot{main result with magnitude}. \slot{Mechanism or welfare implication}.
\end{abstract}
%% END ABSTRACT

%% BEGIN KEYWORDS
\begin{keyword}
\kwd{\slot{Keyword}}
\kwd{\slot{Keyword}}
\kwd{\slot{Keyword}}
\end{keyword}
%% END KEYWORDS

\end{aug}

\end{frontmatter}

%% BEGIN BODY
\section{Introduction}\label{sec:introduction}

A growing literature documents \slot{established fact}. It remains open whether \slot{the paper's question}. This paper answers the question by studying \slot{setting}.

\slot{Why the setting decides the question: the agents, the stakes and the margin the policy moves.}

This paper analyzes the effect of \slot{treatment} by linking \slot{first data source, period} with \slot{second data source, period}. We use \slot{design}, leveraging \slot{the rule, date or assignment that generates the variation}.

We estimate the effect on three outcomes. First, \slot{outcome and estimate with standard error}. Second, \slot{outcome and estimate}. Third, \slot{outcome and estimate}.

These results answer \slot{the question}: \slot{the economic magnitude, incidence or welfare quantity implied by the estimates}.

This paper contributes to three literatures. \slot{Literature one, its benchmark estimate and the estimate this paper adds.} \slot{Literature two.} \slot{Literature three.}

Section~\ref{sec:setting} describes the institutional setting. Section~\ref{sec:data} presents the data. Section~\ref{sec:framework} lays out the conceptual framework. Section~\ref{sec:strategy} describes the empirical strategy. Section~\ref{sec:results} reports the results. Section~\ref{sec:mechanisms} studies the mechanisms. Section~\ref{sec:conclusion} concludes.

\section{Institutional Setting}\label{sec:setting}

\slot{The rule, its history and timing, the agents it reaches and the choice margin it moves.}

\section{Data}\label{sec:data}

\slot{Sources, linkage, unit of observation, period, sample restrictions and outcome construction.} Table~\ref{tab:summary} reports summary statistics.

\begin{table}[t]
\centering
\caption{Summary Statistics}\label{tab:summary}
\begin{threeparttable}
\begin{tabular}{@{}lccc@{}}
\toprule
 & All & Treated & Comparison \\
 & (1) & (2) & (3) \\
\midrule
\slot{Outcome} & 0.000 & 0.000 & 0.000 \\
 & [0.000] & [0.000] & [0.000] \\
\slot{Covariate} & 0.000 & 0.000 & 0.000 \\
 & [0.000] & [0.000] & [0.000] \\
\midrule
\textit{N} & 0 & 0 & 0 \\
\bottomrule
\end{tabular}
\begin{tablenotes}[flushleft]\footnotesize
\item \textit{Notes.} Means and standard deviations (in brackets) for \slot{sample, period and unit}. Column (2) includes \slot{treated units}; column (3) includes \slot{comparison units}.
\end{tablenotes}
\end{threeparttable}
\end{table}

\section{Conceptual Framework}\label{sec:framework}

\slot{Agents, objectives, choice sets, information and equilibrium.} The framework delivers the prediction that Section~\ref{sec:mechanisms} takes to the data.

\begin{assumption}\label{ass:primitive}
\slot{Primitive restriction.}
\end{assumption}

\begin{proposition}\label{prop:main}
Under Assumption~\ref{ass:primitive}, \slot{prediction}.
\end{proposition}

Appendix~\ref{app:proofs} proves Proposition~\ref{prop:main}.

\section{Empirical Strategy}\label{sec:strategy}

We estimate
\begin{equation}\label{eq:main}
y_{it} = \alpha_i + \lambda_{t} + \beta D_{it} + X_{it}'\gamma + \varepsilon_{it},
\end{equation}
where $y_{it}$ is \slot{outcome} for unit $i$ in period $t$, $\alpha_i$ and $\lambda_t$ are unit and period fixed effects, $D_{it}$ indicates \slot{treatment} and $X_{it}$ contains \slot{controls}. The parameter $\beta$ identifies \slot{estimand} under
\begin{equation}\label{eq:identification}
\E\left[\varepsilon_{it}\mid D_{it}, X_{it}, \alpha_i, \lambda_t\right] = 0,
\end{equation}
which states that \slot{the economic content of the condition}. Standard errors are clustered by \slot{level of assignment}.

\section{Results}\label{sec:results}

Figure~\ref{fig:main} presents event-study estimates of the effect of \slot{treatment} on \slot{outcome}. Table~\ref{tab:main} summarizes the estimates.

\begin{figure}[t]
\centering
\IfFileExists{figures/main.pdf}{\includegraphics[width=0.85\textwidth]{figures/main.pdf}}{\fbox{\parbox[c][0.32\textwidth][c]{0.85\textwidth}{\centering figures/main.pdf}}}
\caption{The Effect of \slot{Treatment} on \slot{Outcome}}\label{fig:main}
\begin{minipage}{0.85\textwidth}\footnotesize
The figure presents event-study estimates of the effect of \slot{treatment} on \slot{outcome}. The sample is \slot{sample}. Regressions include \slot{fixed effects}. Error spikes give 95\% confidence intervals, with standard errors clustered by \slot{level}.
\end{minipage}
\end{figure}

\begin{table}[t]
\centering
\caption{Regression Estimates: Effect of \slot{Treatment} on \slot{Outcomes}}\label{tab:main}
\begin{threeparttable}
\begin{tabular}{@{}lcccc@{}}
\toprule
 & \multicolumn{2}{c}{\slot{Outcome A}} & \multicolumn{2}{c}{\slot{Outcome B}} \\
\cmidrule(lr){2-3}\cmidrule(l){4-5}
 & (1) & (2) & (3) & (4) \\
\midrule
\slot{Treatment} & 0.000 & 0.000 & 0.000 & 0.000 \\
 & (0.000) & (0.000) & (0.000) & (0.000) \\
\midrule
Mean of outcome & 0.000 & 0.000 & 0.000 & 0.000 \\
\textit{N} & 0 & 0 & 0 & 0 \\
Unit fixed effects & X & X & X & X \\
Period fixed effects & X & X & X & X \\
Controls &  & X &  & X \\
\bottomrule
\end{tabular}
\begin{tablenotes}[flushleft]\footnotesize
\item \textit{Notes.} Estimates of $\beta$ in equation~\eqref{eq:main} for the outcomes indicated by column headers. Columns (2) and (4) add \slot{controls}. Standard errors (in parentheses) are clustered at the \slot{level}.
\end{tablenotes}
\end{threeparttable}
\end{table}

\section{Mechanisms}\label{sec:mechanisms}

\slot{The test of Proposition~\ref{prop:main} and its estimate; the counterfactual, incidence or welfare calculation.}

\section{Conclusion}\label{sec:conclusion}

\slot{The answer, its magnitude and the decision it informs.}

\begin{appendix}
\section{Proofs}\label{app:proofs}

\begin{proof}[Proof of Proposition~\ref{prop:main}]
\slot{Proof.}
\end{proof}

\section{Additional Results}\label{app:additional}

\slot{Additional tables and figures.}
\end{appendix}
%% END BODY

%% BEGIN BIBLIOGRAPHY
\bibliographystyle{ecta-fullname}
\bibliography{../bibliography}
%% END BIBLIOGRAPHY

\end{document}
